\scanpage{11}
\mspage{(9)}
Ce qui donne à penser que le complexe de DR. p.d. de $\mathcal X_*$ \inserted{à coeff dans $\mathbb Z$ disons} (même dans le cas $k=\mathbb Q$ qu'il contient) donne \unsure{une} information homotopique précise sur $\mathcal X_*$, c'est l'observation qu'en fait, pour un \deleted{module gradué} $k$-module $M$ quelconque, on récupère $M$ à isom canonique près par la connaissance d'un quelconque des tronqués $\tau_q(M\otimes_k S)$ \inserted{\unsure{qu'on appelle|j'appelle} la \uline{$q$ième ombre de $M$}} (quelque grand que soit $\unsure{p|q}\ldots$), et de $=$ l'accouplement $M\otimes N\mathbin{-}P$ est connu quand on connaît, pour $q$ assez grand, l'accouplement correspondant $\tau_q(M_S)\otimes\tau_q(N_S)\mathbin{-}\tau_q(P_S)$.

Donc une \inserted{$k$-}algèbre graduée $H^\bullet$ à degrés $\geq0$ \deleted{est} est connue à isom canonique près quand on connaît la $S$-algèbre bigraduée dont les composantes \inserted{de degré} ``extérieur'' \inserted{$q$} sont les $\tau_q(H^q\otimes_k S)$.... \deleted{Ainsi} Cette ``théorie des ombres'' étant \inserted{\unsure{supposée}} acquise, on voit que la connaissance de $C^{\bullet\bullet}_{\mathrm{DRpd}}(\mathcal X_*,k)$ \inserted{(en tant que $S$-algèbre diff. bigraduée)} implique celle de l'algèbre graduée $H^\bullet(\mathcal X_*,k)$ -- et en raffinant un peu, on \inserted{voit que} $\unsure{F|\vDash}\ R^\bullet\Gamma(\mathcal X_*,k)$ comme objet de la cat. dérivée $D(k)$, \uline{plus} l'accouple-
\mspage{10}
